\chapter{Линейная алгебра}

\section{Характеристический и минимальный многочлен. Кэли--Гамильтон}

\subsection{Теорема Кэли--Гамильтона}
\begin{description}
  \item[Формулировка:] Для любой квадратной матрицы $A$ над коммутативным кольцом: $\chi_A(A)=0$, где $\chi_A(t)=\det(tI-A)$.
  \item[Когда использовать:] Сведение $A^k$ ($k\ge n$) к линейным комбинациям $I,A,\dots,A^{n-1}$. Выражение $A^{-1}$ как многочлена от $A$.
\end{description}

\subsection{Минимальный многочлен $\mu_A$}
\begin{description}
  \item[Формулировка:] Нормированный многочлен наименьшей степени с $\mu_A(A)=0$. Делит $\chi_A$, имеет те же корни. Кратность $\lambda$ в $\mu_A$ равна размеру крупнейшей жордановой клетки для $\lambda$.
  \item[Критерий диагонализуемости:] $A$ диагонализуема над $F$ $\iff$ $\mu_A$ раскладывается над $F$ на различные линейные множители.
\end{description}

\medskip

\section{След (trace) и спектр произведения}

\subsection{Свойства следа}
$\operatorname{tr}(AB)=\operatorname{tr}(BA)$, $\operatorname{tr}(A)=\sum\lambda_i$, $\operatorname{tr}(A^k)=\sum\lambda_i^k$. След инвариантен относительно подобия.

\subsection{Спектр $AB$ vs $BA$}
$A\in M_{m\times n}$, $B\in M_{n\times m}$: $\det(\lambda I_n-BA)=\lambda^{n-m}\det(\lambda I_m-AB)$. Ненулевые собственные значения $AB$ и $BA$ совпадают. При $m=n$: $\chi_{AB}=\chi_{BA}$. Следствие: $\det(I+AB)=\det(I+BA)$.

\subsection{Тождества Ньютона}
$p_k=\operatorname{tr}(A^k)$, $\chi_A(t)=\sum(-1)^k e_k t^{n-k}$. Тогда
$$p_k - e_1 p_{k-1} + e_2 p_{k-2} - \dots + (-1)^{k-1}k e_k = 0,\quad k=1,\dots,n.$$
Над полем характеристики $0$ знание $p_1,\dots,p_n$ восстанавливает $\chi_A$.

\subsection{Критерий нильпотентности}
Над полем характеристики $0$: $A$ нильпотентна $\iff$ $\operatorname{tr}(A^k)=0$ для $k=1,\dots,n$.

\medskip

\section{Ранг}

\subsection{Неравенства Сильвестра и Фробениуса}
Сильвестр: $\operatorname{rank}(A)+\operatorname{rank}(B)-n\le\operatorname{rank}(AB)\le\min(\operatorname{rank}(A),\operatorname{rank}(B))$.
Фробениус: $\operatorname{rank}(AB)+\operatorname{rank}(BC)\le\operatorname{rank}(B)+\operatorname{rank}(ABC)$.
Следствие: $AB=0\Rightarrow\operatorname{rank}(A)+\operatorname{rank}(B)\le n$.

\medskip

\section{Жорданова форма и спектральная теорема}

\subsection{Жорданова нормальная форма}
Над $\mathbb{C}$: $A=SJS^{-1}$, $J=\bigoplus J_{m_k}(\lambda_k)$, $J_m(\lambda)=\lambda I_m+N_m$. Число блоков размера $\ge k$ для $\lambda$ равно $\dim\ker(A-\lambda I)^k-\dim\ker(A-\lambda I)^{k-1}$.

\subsection{Спектральная теорема для нормальных матриц}
$A\in M_n(\mathbb{C})$ нормальна ($AA^*=A^*A$) $\iff$ существует унитарная $U$ с $U^*AU=\operatorname{diag}(\lambda_1,\dots,\lambda_n)$.
\begin{itemize}
  \item Эрмитова ($A=A^*$): $\lambda_i\in\mathbb{R}$.
  \item Унитарная ($A^*A=I$): $|\lambda_i|=1$.
  \item Вещественная симметричная: ортогонально диагонализуема над $\mathbb{R}$.
\end{itemize}

\subsection{Триангуляризация Шура}
Над $\mathbb{C}$: существует унитарная $U$ и верхнетреугольная $T$ с $U^*AU=T$, $T_{ii}=\lambda_i$.

\medskip

\section{Положительно (полу)определённые матрицы}

\subsection{Эквивалентные критерии PSD}
Для эрмитовой $A$ эквивалентны:
\begin{enumerate}
  \item $x^*Ax\ge0$ для всех $x$;
  \item $\lambda_i\ge0$;
  \item $A=B^*B$ для некоторой $B$;
  \item все главные миноры $\ge0$.
\end{enumerate}
Для PD: строгие неравенства; в (4) достаточно ведущих главных миноров (критерий Сильвестра).

\subsection{Представление Грама}
$A\succeq0$ $\iff$ $A=B^*B$ $\iff$ $A_{ij}=\langle v_i,v_j\rangle$. $\operatorname{rank}(A)=\dim\operatorname{span}(v_i)$. Любая PSD-матрица --- матрица Грама некоторого набора векторов.

\subsection{Разложение Холецкого}
$A\succ0$ $\iff$ существует единственная нижнетреугольная $L$ с $L_{ii}>0$ и $A=LL^*$.

\subsection{Квадратный корень}
Для $A\succeq0$ существует единственный $A^{1/2}\succeq0$ с $(A^{1/2})^2=A$ (через спектр: $A^{1/2}=U\Lambda^{1/2}U^*$).

\subsection{Дополнение Шура}
Для $M=\begin{pmatrix}A&B\\ B^*&D\end{pmatrix}$ с $D\succ0$: $M\succ0$ $\iff$ $A-BD^{-1}B^*\succ0$.

\subsection{Порядок Лёвнера}
$A\succeq B$ $\iff$ $A-B\succeq0$. Свойства: $A\succeq B\Rightarrow C^*AC\succeq C^*BC$; $0\prec A\preceq B\Rightarrow B^{-1}\preceq A^{-1}$.

\subsection{Сингулярное разложение (SVD)}
Любая $A\in\mathbb{C}^{m\times n}$: $A=U\Sigma V^*$, $U,V$ унитарны, $\Sigma$ диагональна с $\sigma_1\ge\cdots\ge\sigma_r>0$. $\sigma_i=\sqrt{\lambda_i(A^*A)}$, $r=\operatorname{rank}(A)$.

\subsection{Полярное разложение}
Любая $A$: $A=UP$, $U$ унитарна, $P=(A^*A)^{1/2}\succeq0$.

\medskip

\section{Евклидовы пространства и матрица Грама}

\subsection{Матрица Грама}
$G_{ij}=\langle v_i,v_j\rangle$. $G$ симметрична и PSD. $G\succ0$ $\iff$ $v_i$ линейно независимы. $\det G=(\operatorname{Vol})^2$.

\subsection{Неравенство Адамара}
$|\det A|\le\prod\|a_i\|$ (столбцы). Равенство $\iff$ столбцы ортогональны. Для PSD: $\det H\le\prod H_{ii}$.

\subsection{Ортогональная проекция}
Для $W\subseteq V$: $P_W^2=P_W$, $P_W^T=P_W$, $\operatorname{im}P_W=W$, $\ker P_W=W^\perp$. $\|x-P_Wx\|=\min_{y\in W}\|x-y\|$. Если столбцы $A$ --- базис $W$: $P_W=A(A^TA)^{-1}A^T$.

\subsection{Неравенство Бесселя}
Для ортонормированной системы $e_1,\dots,e_k$: $\sum_{i=1}^k|\langle x,e_i\rangle|^2\le\|x\|^2$. Равенство $\iff$ $x\in\operatorname{span}(e_i)$.

\subsection{Неравенство Коши--Шварца}
$|\langle u,v\rangle|\le\|u\|\cdot\|v\|$, равенство $\iff$ $u,v$ линейно зависимы.

\medskip

\section{Аналитическая геометрия}

\subsection{Прямая на плоскости}
\begin{itemize}
\item \textit{Параметрически:} $(x,y) = (x_0,y_0) + t(a,b)$.
\item \textit{Общее уравнение:} $Ax+By+C=0$, нормаль $(A,B)$.
\item \textit{Расстояние от точки $(x_0,y_0)$:} $\dfrac{|Ax_0+By_0+C|}{\sqrt{A^2+B^2}}$.
\item \textit{Угол между прямыми:} $\cos\theta = \dfrac{|A_1A_2+B_1B_2|}{\sqrt{A_1^2+B_1^2}\sqrt{A_2^2+B_2^2}}$.
\end{itemize}

\subsection{Плоскость и прямая в пространстве}
\begin{itemize}
\item \textit{Плоскость:} $Ax+By+Cz+D=0$, нормаль $(A,B,C)$.
\item \textit{Расстояние от точки до плоскости:} $\dfrac{|Ax_0+By_0+Cz_0+D|}{\sqrt{A^2+B^2+C^2}}$.
\item \textit{Прямая в $\mathbb{R}^3$:} $(x,y,z) = (x_0,y_0,z_0) + t(a,b,c)$.
\item \textit{Расстояние между скрещивающимися прямыми:} $d = \dfrac{|(\mathbf{r}_2-\mathbf{r}_1)\cdot(\mathbf{a}_1\times\mathbf{a}_2)|}{\|\mathbf{a}_1\times\mathbf{a}_2\|}$.
\end{itemize}

\subsection{Конические сечения (коники)}

Общее уравнение второй степени: $Ax^2+Bxy+Cy^2+Dx+Ey+F=0$.

Тип кривой определяется дискриминантом $\Delta = B^2-4AC$:
\begin{itemize}
\item $\Delta<0$: эллипс (в частном случае окружность при $A=C$, $B=0$).
\item $\Delta=0$: парабола.
\item $\Delta>0$: гипербола.
\end{itemize}

\subsubsection{Эллипс}

$\dfrac{x^2}{a^2}+\dfrac{y^2}{b^2}=1$, $a\ge b>0$.
\begin{itemize}
\item Фокусы $(\pm c,0)$, $c=\sqrt{a^2-b^2}$.
\item $r_1+r_2 = 2a$ (сумма расстояний до фокусов постоянна).
\item Эксцентриситет $e=c/a<1$.
\item Директрисы $x=\pm a/e$.
\item Площадь $S=\pi ab$.
\end{itemize}

\subsubsection{Гипербола}

$\dfrac{x^2}{a^2}-\dfrac{y^2}{b^2}=1$.
\begin{itemize}
\item Фокусы $(\pm c,0)$, $c=\sqrt{a^2+b^2}$.
\item $|r_1-r_2| = 2a$ (модуль разности расстояний до фокусов постоянен).
\item Эксцентриситет $e=c/a>1$.
\item Директрисы $x=\pm a/e$.
\item Асимптоты $y=\pm(b/a)x$.
\end{itemize}

\subsubsection{Парабола}

$y^2 = 2px$, $p>0$.
\begin{itemize}
\item Фокус $(p/2,0)$.
\item Директриса $x=-p/2$.
\item $r = d$ (расстояние до фокуса равно расстоянию до директрисы).
\item Эксцентриситет $e=1$.
\end{itemize}

\medskip

\section{Билинейные и квадратичные формы. Сигнатура}

\subsection{Закон инерции Сильвестра}
Для симметричной $A\in\mathbb{R}^{n\times n}$ сигнатура $(p,q,r)$ --- числа положительных, отрицательных и нулевых собственных значений --- инвариантна относительно конгруэнтности $A\mapsto C^TAC$.

\subsection{Формула Якоби}
Если все ведущие миноры $\Delta_k\neq0$, то существует верхнетреугольная $U$ с $U^TAU=\operatorname{diag}(\Delta_1,\Delta_2/\Delta_1,\dots,\Delta_n/\Delta_{n-1})$.

\subsection{Одновременная диагонализация}
Если $A\succ0$ и $B$ симметрична, то существует $C$ с $C^TAC=I$, $C^TBC=\operatorname{diag}(\mu_i)$. Числа $\mu_i$ --- корни $\det(B-\mu A)=0$, все вещественны.

\medskip

\section{Определители}

\subsection{Вандермонд}
$\det(x_j^{i-1})=\prod_{i<j}(x_j-x_i)$. Обобщение: для монических многочленов $P_{i-1}$ степени $i-1$ определитель тот же.

\subsection{Блочные формулы}
$\det\begin{pmatrix}A&B\\ C&D\end{pmatrix}=\det A\cdot\det(D-CA^{-1}B)$ (при обратимой $A$). Если $CD=DC$, то $\det\begin{pmatrix}A&B\\ C&D\end{pmatrix}=\det(AD-BC)$.

\subsection{Коши--Бине}
Для $A\in M_{m\times n}$, $B\in M_{n\times m}$, $m\le n$:
$$\det(AB)=\sum_{S,\,|S|=m}\det(A_{[m],S})\cdot\det(B_{S,[m]}).$$
Следствие: $\det(A^TA)=\sum_S\det(A_{S})^2\ge0$.

\subsection{Тождество Сильвестра}
$\det(I_m+AB)=\det(I_n+BA)$. Частный случай: $\det(I+uv^T)=1+v^Tu$.

\subsection{Возмущение ранга один}
$\det(A+uv^T)=\det(A)\cdot(1+v^TA^{-1}u)$. Для ранга $k$: $\det(A+UV^T)=\det(A)\cdot\det(I_k+V^TA^{-1}U)$.

\subsection{Определитель Коши}
$\det(1/(a_i+b_j))=\dfrac{\prod_{i<j}(a_i-a_j)(b_i-b_j)}{\prod_{i,j}(a_i+b_j)}$.

\subsection{Циркулянт}
$C=\operatorname{circ}(c_0,\dots,c_{n-1})$: $\det C=\prod_{k=0}^{n-1} p(\omega^k)$, где $p(x)=\sum c_k x^k$, $\omega=e^{2\pi i/n}$.

\medskip

\section{Специальные классы матриц}

\subsection{Нильпотентная матрица}
$N^k=0$. Эквивалентно: $\sigma(N)=\{0\}$, $\chi_N(t)=t^n$. Жорданова форма --- блоки $J_{m_i}(0)$. $(I-N)^{-1}=I+N+\cdots+N^{n-1}$.

\subsection{Идемпотент}
$P^2=P$. Тогда $P$ диагонализуема, $V=\operatorname{im}P\oplus\ker P$, $\sigma(P)\subseteq\{0,1\}$. $\operatorname{rank}(P)=\operatorname{tr}(P)$. $I-P$ --- тоже идемпотент.

\subsection{Инволюция}
$A^2=I$, char $\neq2$. Тогда $\sigma(A)\subseteq\{\pm1\}$, $A$ диагонализуема. $V=V_+\oplus V_-$, $V_\pm=\ker(A\mp I)$. Проекторы $P_\pm=(I\pm A)/2$.

\subsection{Коммутирующие матрицы}
Если $AB=BA$, то: существует общий собственный вектор; семейство попарно коммутирующих матриц одновременно триангуляризуемо; если все диагонализуемы --- одновременно диагонализуемы.

\subsection{Уравнение Сильвестра}
$AX-XB=C$ имеет единственное решение $\iff$ $\sigma(A)\cap\sigma(B)=\emptyset$.

\subsection{Коммутатор}
$\operatorname{tr}(AB-BA)=0$. При char$=0$ уравнение $AB-BA=I$ неразрешимо.

\medskip

\section{Произведение Кронекера}

$(A\otimes B)(C\otimes D)=(AC)\otimes(BD)$. $(A\otimes B)^T=A^T\otimes B^T$, $(A\otimes B)^{-1}=A^{-1}\otimes B^{-1}$.

$\sigma(A\otimes B)=\{\lambda_i\mu_j\}$. $\operatorname{tr}(A\otimes B)=\operatorname{tr}(A)\operatorname{tr}(B)$. $\det(A\otimes B)=(\det A)^m(\det B)^n$.

$\operatorname{vec}(AXB)=(B^T\otimes A)\operatorname{vec}(X)$.

$A\oplus B=A\otimes I_m+I_n\otimes B$ (кронекерова сумма). $\sigma(A\oplus B)=\{\lambda_i+\mu_j\}$.

\medskip

\section{Матрицы над $\mathbb{F}_p$}

\subsection{Тождество Фробениуса}
Для коммутирующих $A,B$ над полем char $p$: $(A+B)^{p^k}=A^{p^k}+B^{p^k}$.

\subsection{Порядок $\operatorname{GL}_n(\mathbb{F}_q)$}
$|\operatorname{GL}_n(\mathbb{F}_q)|=\prod_{k=0}^{n-1}(q^n-q^k)$. Число $k$-мерных подпространств: $\binom{n}{k}_q$.

\subsection{Теорема Шевалле--Уорнинга}
$f_1,\dots,f_r\in\mathbb{F}_q[x_1,\dots,x_n]$, $\sum\deg f_i<n$. Число общих нулей в $\mathbb{F}_q^n$ делится на $p$. Если есть тривиальный ноль --- есть и нетривиальный.

\subsection{Комбинаторный Nullstellensatz Алона}
$f\in F[x_1,\dots,x_n]$, $\deg f=\sum t_i$, коэффициент при $\prod x_i^{t_i}$ ненулевой. Для $S_i\subseteq F$ с $|S_i|>t_i$ существует $s\in\prod S_i$ с $f(s)\neq0$.

\medskip

\section{Матричная экспонента}

\subsection{Вычисление}
$e^{tJ_m(\lambda)}=e^{t\lambda}\sum_{k=0}^{m-1}\frac{t^k}{k!}N_m^k$. $e^{tA}=S\,e^{tJ}\,S^{-1}$.

$\sigma(e^A)=e^{\sigma(A)}$, $\det e^A=e^{\operatorname{tr}A}$.

$AB=BA$ $\Rightarrow$ $e^{A+B}=e^A e^B$.

\subsection{Оценки}
$\|A^n\|\le C\cdot n^{m-1}\rho(A)^n$, $\|e^{tA}\|\le C(1+t^{m-1})e^{\alpha(A)t}$ ($t\ge0$). Формула Гельфанда: $\rho(A)=\lim\|A^n\|^{1/n}$.

\medskip

\section{Целочисленные матрицы (Smith normal form)}

\subsection{Нормальная форма Смита}
Для $A\in M_{m\times n}(\mathbb{Z})$ ранга $r$ существуют унимодулярные $U,V$ с $UAV=\operatorname{diag}(d_1,\dots,d_r,0,\dots)$, $d_1\mid d_2\mid\cdots\mid d_r$. $\prod d_i=|\det A|$ (для квадратной).

\subsection{Инвариантные факторы через миноры}
$\Delta_k(A)=\gcd$ всех $k\times k$ миноров $A$. $d_k=\Delta_k/\Delta_{k-1}$.

\subsection{Структура коядра}
$\mathbb{Z}^n/A\mathbb{Z}^n\cong\bigoplus\mathbb{Z}/d_i\oplus\mathbb{Z}^{n-r}$. Для невырожденной $A$: $|\mathbb{Z}^n/A\mathbb{Z}^n|=|\det A|$.

\medskip

\section{Минимаксные характеризации и неравенства}

\subsection{Куранта--Фишера}
Для эрмитовой $A$ с $\lambda_1\ge\cdots\ge\lambda_n$:
$$\lambda_k=\max_{\dim V=k}\min_{x\in V,\|x\|=1}x^*Ax=\min_{\dim V=n-k+1}\max_{x\in V,\|x\|=1}x^*Ax.$$
Для сингулярных чисел: $\sigma_k(A)=\max_{\dim V=k}\min_{x\in V,\|x\|=1}\|Ax\|$.

\subsection{Перемежение Коши}
Для главной подматрицы $B$ размера $n-1$: $\lambda_1\ge\mu_1\ge\lambda_2\ge\mu_2\ge\cdots\ge\mu_{n-1}\ge\lambda_n$.

\subsection{Неравенства Вейля}
Для эрмитовых $A,B$: $\lambda_{i+j-1}(A+B)\le\lambda_i(A)+\lambda_j(B)$.

\subsection{Неравенство Шура}
$\sum|\lambda_i|^2\le\|A\|_F^2$. Равенство $\iff$ $A$ нормальна.

\subsection{Неравенство Минковского}
Для $A,B\succeq0$: $\det(A+B)^{1/n}\ge\det(A)^{1/n}+\det(B)^{1/n}$.
